%
%
%
%
%
%
\chapter{Background}
\label{cha:background}
\noindent
This chapter establishes the theoretical foundation necessary for the thesis.
We begin by specifying the \While language, from \emph{Semantics with Applications: An Appetizer}
which is our primary study which throughout this work. This provides the concrete syntax required to
explore the formal semantics of programming languages—the mathematical bridge
between a program's notation and its behavior. We then delve into Operational
Semantics, detailing both the Natural and Structural
approaches, followed by the primary proof techniques used in their analysis.
Finally, we introduce the Lean 4 proof assistant, detailing the mechanisms used
to translate these textbook definitions into machine-checked proofs~\citep{theorem_proving_2026}.

%
%
%
\section{The \While Language Framework}
\label{background:sec:while}
The formal language analyzed throughout this work is \While, an imperative, deterministic toy language defined by \citet{nielson_semantics_2007}. To read the formal proofs and translations compared in this thesis, a basic understanding of the language's structural notation is required.

\subsection*{Syntactic Structure and BNF Notation}
The abstract syntax of \While is defined inductively over five core syntactic categories, each associated with a specific meta-variable used in proofs:
\begin{itemize}
    \item $n \in \Num$ (Numerals / digits)
    \item $x \in \Var$ (Variable names)
    \item $a \in \Aexp$ (Arithmetic expressions, e.g., $x + 1$)
    \item $b \in \Bexp$ (Boolean expressions, e.g., $a_1 \leq a_2$)
    \item $S \in \Stm$ (Statements / commands, e.g., loops and assignments)
\end{itemize}
The grammar of these categories is specified using Backus-Naur Form (BNF) notation. In this notation, the symbol \mbox{\(::=\)} means \emph{"can be structured as"}, and the vertical bar \mbox{\(\mid\)} denotes an \emph{"alternative choice"}. For example, the statement grammar is written as:
\[
S ::= x := a \mid \mathbf{skip} \mid S_1; S_2 \mid \mathbf{if}\ b\ \mathbf{then}\ S_1\ \mathbf{else}\ S_2 \mid \mathbf{while}\ b\ \mathbf{do}\ S
\]
This indicates that a statement $S$ is either an assignment, a no-op (\textbf{skip}), a sequential composition of two statements separated by a semicolon, a conditional branch, or a loop.

\subsection*{Semantic States and Denotational Notation}
To give the syntax mathematical meaning, we map expressions to values using explicit environments and evaluation functions:
\begin{itemize}
    \item \textbf{Program State (\State):} Represented by the meta-variable $s$, a state is a total function mapping variables to integers (\(s : \mathbf{Var} \to \mathbb{Z}\)). The notation \(s \, x\) denotes looking up the integer value of variable $x$ in state $s$.
    \item \textbf{State Override Notation \(s[x \mapsto v]\):} This represents a lookup function identical to $s$, except that the variable $x$ is explicitly updated to map to the new integer value $v$.
    \item \textbf{Semantic Brackets \([\![\cdot]\!]\):} Double brackets (Oxford brackets) are used to separate the syntax of the programming language from the mathematical meta-language evaluating it. 
\end{itemize}
The textbook defines three compositional semantic functions using this bracket notation:
\begin{table}[H]
\centering
\begin{tabularx}{\textwidth}{llX}
\hline
\textbf{Function Type Signature} & \textbf{Notation Example} & \textbf{Meaning} \\ \hline
$\mathcal{N} : \mathbf{Num} \to \mathbb{Z}$ & $\mathcal{N}[\![1]\!]$ & The syntactic digit $1$ becomes the integer $1$. \\[1.8em]
$\mathcal{A} : \mathbf{Aexp} \to (\mathbf{State} \to \mathbb{Z})$ & $\mathcal{A}[\![a]\!]s$ & Evaluates arithmetic expression $a$ inside state $s$. \\[1.8em]
$\mathcal{B} : \mathbf{Bexp} \to (\mathbf{State} \to \{\mathbf{tt}, \mathbf{ff}\})$ & $\mathcal{B}[\![b]\!]s$ & Evaluates boolean expression $b$ to true (\textbf{tt}) or false (\textbf{ff}) in state $s$. \\ \hline
\end{tabularx}
\end{table}


\section{Operational Semantics Frameworks}
\label{background:sec:operational_semantics}
Operational semantics describes how a program executes by tracking changes to an abstract machine. This machine is represented as a \textbf{configuration}, written inside angle brackets \(\langle S, s \rangle\), which pairs a statement $S$ yet to be executed with the current program state $s$. 

This thesis contrasts the translation of two operational paradigms defined by \citet{nielson_semantics_2007}: Natural Semantics and Structural Operational Semantics.

\subsection*{Natural Semantics (Big-Step Notation)}
Natural Semantics (NS) describes the execution of a statement running completely to its conclusion. The transition relation is denoted by a single arrow (\(\to\)):
\[
\langle S, s \rangle \to s'
\]
This is read as: \emph{"The execution of statement $S$ starting in state $s$ terminates completely, resulting in the final state $s'$."} Rules in this framework are written as fractions (inference rules), where the statements above the horizontal line are the required premises, and the statement below is the conclusion.

\subsection*{Structural Operational Semantics (Small-Step Notation)}
Structural Operational Semantics (SOS) models execution at a higher resolution, tracking individual, atomic steps of computation. The transition relation is denoted by a double arrow (\(\Rightarrow\)). SOS distinguishes between two types of targets:
\begin{itemize}
    \item \textbf{Intermediate Transition:} $\langle S, s \rangle \Rightarrow \langle S', s' \rangle$ \\
    Read as: \emph{"Executing the first step of statement $S$ in state $s$ modifies the state to $s'$, leaving statement $S'$ remaining to be executed."}
    \item \textbf{Terminal Transition:} $\langle S, s \rangle \Rightarrow s'$ \\
    Read as: \emph{"Executing the final step of statement $S$ in state $s$ finishes execution and results in the terminal state $s'$."}
\end{itemize}
To represent multi-step program execution from an initial state to a final state, SOS uses the reflexive transitive closure notation, written with an asterisk: \(\Rightarrow^*\). The notation \(\langle S, s \rangle \Rightarrow^* s'\) means the configuration reaches state $s'$ after zero or more sequential small steps.


%
%
%
\section{The Lean 4 Proof Assistant}
\label{background:sec:lean_4}
Lean 4 is a functional programming language and an interactive theorem prover
based on \textit{Dependent Type Theory} \citep{theorem_proving_2026}.
Lean can be viewed as a rigorous environment for defining mathematical objects
and constructing formal proofs that are mechanically verified by a small logical so-called kernel.

\subsection{Defining Syntax: Inductive Types}
In the textbook \emph{Semantics with Applications: An Appetizer} by \cite{nielson_semantics_2007},
the syntax of a programming language is often defined using Backus-Naur Form (BNF).
In Lean, we represent these syntactic categories using \texttt{inductive} types.
An inductive type defines a set by listing its constructors.
For example, the syntactic category of arithmetic expressions (\textit{Aexp}) is defined as:

\begin{Lean_Code}
inductive Aexp where
  | num (n : Num)          -- Constants
  | var (x : Var)          -- Variables
  | add (a1 a2 : Aexp)     -- Addition
\end{Lean_Code}
\noindent
Where the green text starting with \texttt{--} represents a comment.
This maps directly to the mathematical notion of a set defined by structural induction.
Each constructor (e.g., \texttt{add}) is a function
that takes existing elements of the set and produces a new one.

\subsection{Defining Semantics: Relations as Propositions}
The evaluation of a statement is written as a relation: $\langle S, s \rangle \to s'$. 
In Lean, this relation is defined as an \texttt{inductive} type that returns
a \texttt{Prop} (Proposition).
This is a crucial distinction: \textbf{Data Types} (like \texttt{Aexp}) are used to build programs.
Whilst \textbf{Propositions} (like \texttt{big\_step}) are used to make claims about those programs.
Due to propositions being something  that is either true or false, whilst data types are just syntactic sugar.

\noindent
When we define \texttt{big\_step} inductively, each constructor represents an
\textbf{inference rule}.
We can then combine this with custom notation so that the original Lean code of:

\begin{Lean_Code}
    big_step S s s'
\end{Lean_Code}
can be transformed into
\begin{Lean_Code}
    §\(\langle\)§S, s§\(\rangle\)§ §\(\to_{ns}\)§ s'
\end{Lean_Code}
by using the following
\begin{Lean_Code}
notation "§\color{purple}\(\langle\)§" S ", " s "§\color{purple}\(\rangle\)§ §\color{purple}\(\to_{ns}\)§ " s' => big_step S s s'
\end{Lean_Code}
For instance, the rule for statement composition (\texttt{comp}) in Lean is:

\begin{Lean_Code}
| comp (h_left : §\(\langle\)§S§\(_1\)§, s§\(\rangle\)§ §\(\to_{ns}\)§ s') (h_right : §\(\langle\)§S§\(_2\)§, s'§\(\rangle\)§ §\(\to_{ns}\)§ s'') :
    §\(\langle\)§S§\(_1\)§; S§\(_2\)§, s§\(\rangle\)§ §\(\to_{ns}\)§ s''
\end{Lean_Code}
This code defines a rule of inference. To a reader familiar with natural deduction from a course in mathematical logic,
this is a direct mapping of the standard horizontal-bar notation:
\[
\frac{\langle S_1, s \rangle \to_{ns} s' \quad \langle S_2, s' \rangle \to_{ns} s''}{\langle S_1; S_2, s \rangle \to_{ns} s''}
\]
The names \texttt{h\_left} and \texttt{h\_right} represent the \textit{hypotheses} or \textit{premises}
required to reach the conclusion below the colon.

\subsection{The Proof State and Tactics}
Unlike a prose proof, which is a static piece of text,
a proof in Lean is constructed using \emph{tactics}.
A tactic is a command that transforms a "goal" (the statement to be proven) into simpler sub-goals.

When reading the proofs in this thesis,
it is helpful to understand the workflow:
Lean maintains a \textit{proof state} consisting of the current \emph{context}
(known facts and variables) and the \emph{goal}.
A tactic like \texttt{intro} moves an assumption into the context,
while \texttt{cases} performs a proof by exhaustion,
forcing Lean to consider every possible way a specific relation could have been formed.

\subsection{Mapping Lean to Mathematical Logic}
The following table provides a quick reference for mapping Lean's core keywords
to their respective functionalities in standard mathematical and logical discourse.

\begin{figure}[H]
\centering
\begin{tabularx}{\textwidth}{lXX}
\hline
\textbf{Lean Keyword} & \textbf{Logical Equivalent} & \textbf{Functionality} \\ \hline
\texttt{inductive}    & Inductive Definition      & Defines a set or a relation via a list of rules/constructors. \\ \\
\texttt{Prop}         & Proposition               & The type of statements that have truth values. \\ \\
\texttt{theorem}      & Theorem / Lemma           & A claim that requires a formal proof term. \\ \\
\texttt{intro}        & Implication Introduction  & "Assume $P$ is true" or "Let $x$ be an arbitrary element." \\ \\
\texttt{exact}        & Q.E.D.                    & Completes a goal by providing the required evidence directly. \\ \\
\texttt{cases}        & Proof by Cases       & Deconstructs an object or relation into its possible forms. \\ \\
\texttt{apply}        & Rule application (implication elimination) & Uses a known rule or theorem to transform or discharge the current goal. \\ \\
\texttt{rfl}          & Reflexivity               & Proves an equality of the form $A = A$ by definition. \\ \hline
\end{tabularx}
\caption{Mapping of Lean 4 keywords to mathematical concepts.}
\end{figure}

\subsection{Notation Conventions}
To bridge the gap between code and literature,
it is possible to define custom notation and operations as demonstrated earlier.
This allows us to mimic notation used by the pen-and-paper proofs.
This can allow for a custom Number type to use the addition sign "+"
to mark an addition of two numbers, instead of using the full inductive name such-as 
"\texttt{Number.Add N1 N2}", and instead do "\texttt{N1 + N2}".
Notations can be defined by using the keyword "\texttt{notation}",
and operations can be defined using "\texttt{infixl}" and "\texttt{infixr}"

